\clearpage\fancyhead[L]{\small COMET.}\hypertarget{native-340-lower}{}\label{native-340-lower}
\begin{center}\textbf{COMET.}\par\vspace{4pt}\hrule\vspace{4pt}\fontsize{10}{12}\selectfont Monatliche Correspondenz. Vol. XXIV. pp. 304\,..\,307. September 1811.\par\vspace{4pt}\hrule\vspace{4pt}\end{center}

\par\indent ---\;---\;--- Professor \textsc{Gauss} shared with us some interesting results of observation and calculation. We give the relevant passages from several of his letters. ---\;---\;--- Göttingen, 8 Sept. “I first saw the comet on 22 August, but the city severely limits the view north from the observatory, so I could not observe it until 3 September. I have tried to determine its position from angular distances measured with the sextant. With the fine stand you procured for us through the skilled instrument maker \textsc{Körner} in Weimar, this is a very convenient way to observe and seems comparatively very accurate as well. I have yet to reduce the observations of 3 September. For comparison on each of the following days I used $\alpha$ Aurigae and $\alpha$ Lyrae. I reduced those observations with great care and found the following results.
\Needspace{9\baselineskip}
\begin{center}\fontsize{10}{12}\selectfont\setlength{\tabcolsep}{4pt}\renewcommand{\arraystretch}{1.15}\begin{tabular}{r|c|c|c}
1811 & \shortstack{Mean time\\in Göttingen} & Comet R.A. & \shortstack{Northern\\declination}\\\hline
Sept. 4 & $8^{\mathrm u}\,28^{\mathrm m}\,47^{\mathrm s}$ & $158^\circ\,25'\,24''$ & $39^\circ\,18'\,2''$\\
6 & $8\;48\;38$ & $160\;23\;16$ & $40\;14\;16$\\
7 & $8\;57\;6$ & $161\;26\;10$ & $40\;41\;54$\\
\end{tabular}\end{center}

\SourceLeaf{341}\par\indent Using these observations and nine by Mr. von \textsc{Zach}, I improved my parabolic elements\footnote{Monatl. Corresp., August 1811 issue, p. 180. [p. 326 of this volume.]} and obtained the following.
\Needspace{10\baselineskip}
\begin{center}\fontsize{10}{12}\selectfont\setlength{\tabcolsep}{4pt}\renewcommand{\arraystretch}{1.15}\begin{tabular}{l}
Time of perihelion 1811 Sept. 12\\
\quad $5^{\mathrm u}\,21^{\mathrm m}\,15^{\mathrm s}$ Mean time in Göttingen\\
\end{tabular}\end{center}
\begin{center}\fontsize{10}{12}\selectfont\setlength{\tabcolsep}{4pt}\renewcommand{\arraystretch}{1.15}\begin{tabular}{llr}
Logarithm of the least distance & .\;.\;.\;.\;.\;.\;.\;. & $0.017060$\\
Longitude of perihelion & .\;.\;.\;.\;.\;.\;.\;.\;.\;.\;.\;.\;. & $75^\circ\,17'\,34''$\\
Longitude of the ascending node & .\;.\;.\;.\;.\;.\;.\;.\;. & $140\;24\;13$\\
Inclination of the orbit & .\;.\;.\;.\;.\;.\;.\;.\;.\;.\;.\;.\;.\;. & $73\;7\;16$\\
\end{tabular}\end{center}

\par\indent There still seems to be no definite sign that the orbit is elliptical. I have only taken the time to look at the comet through a large telescope a few times so far. On 7 September the shape of its tail was very remarkable. It curved away from the comet in two branches. The two branches did not start at the comet itself but joined a short distance from it, separated from it by a dark gap, so that they enclosed the comet as a parabola encloses its focus. In more than one respect this comet is certainly one of the most remarkable ever observed. Professor \textsc{Harding} and I could see no trace of a nucleus as such in the comet.” ---\quad Göttingen, 16 Sept. Since you are now observing the comet on the meridian, I will end my observations this evening. Here is what I have obtained since my last letter.
\Needspace{10\baselineskip}
\begin{center}\fontsize{10}{12}\selectfont\setlength{\tabcolsep}{4pt}\renewcommand{\arraystretch}{1.15}\begin{tabular}{r|c|c|c}
 & \shortstack{Mean time\\in Göttingen} & Comet R.A. & \shortstack{Northern\\declination}\\\hline
Sept. 9 & $8^{\mathrm u}\,37^{\mathrm m}\,53^{\mathrm s}$ & $163^\circ\,35'\,5''$ & $41^\circ\,39'\,13''$\\
14 & $8\;45\;5$ & $169\;44\;58$ & .\quad.\quad.\quad.\\
15 & $7\;39\;57$ & $171\;3\;58$ & $44\;24\;57$\\
16 & $7\;57\;16$ & $172\;31\;4$ & $44\;51\;23$\\
\end{tabular}\end{center}

\Needspace{24\baselineskip}\par\indent Using my improved elements, I have calculated the following ephemeris for the comet’s motion through the end of December.
\begin{center}\fontsize{10}{12}\selectfont\setlength{\tabcolsep}{3pt}\renewcommand{\arraystretch}{1.15}\begin{tabular}{rr|c|c|c|c}
\multicolumn{2}{c|}{\shortstack{Mean time\\in Göttingen}} & \shortstack{Right\\ascension} & \shortstack{Northern\\declination} & \shortstack{Log. of the comet’s\\distance from Earth} & \shortstack{Brightness}\\\hline
1811 Sept. 12 & $5^{\mathrm u}$ & $166^\circ\,54'$ & $42^\circ\,58'$ & $0.2079$ & $0.355$\\
17 & $0$ & $173\;46$ & $45\;15$ & $0.1836$ & $0.394$\\
22 & $6$ & $182\;2$ & $47\;17$ & $0.1594$ & $0.432$\\
.27 & $6$ & $192\;19$ & $48\;48$ & $0.1366$ & $0.465$\\
Oct. 2 & $7$ & $203\;24$ & $49\;34$ & $0.1160$ & $0.490$\\
7 & $7$ & $215\;58$ & $49\;7$ & $0.1004$ & $0.501$\\
12 & $8$ & $228\;45$ & $47\;16$ & $0.0911$ & $0.492$\\
22 & $9$ & $251\;23$ & $39\;46$ & $0.0963$ & $0.421$\\
Nov. 1 & $10$ & $267\;58$ & $30\;32$ & $0.1312$ & $0.310$\\
11 & $11$ & $279\;38$ & $22\;18$ & $0.1851$ & $0.208$\\
21 & $12$ & $288\;13$ & $14\;55$ & $0.2415$ & $0.138$\\
Dec. 1 & $13$ & $294\;54$ & $9\;42$ & $0.2976$ & $0.093$\\
11 & $14$ & $300\;24$ & $5\;49$ & $0.3487$ & $0.064$\\
21 & $15$ & $305\;7$ & $2\;55$ & $0.3938$ & $0.046$\\
31 & $16$ & $309\;16$ & $0\;48$ & $0.4330$ & $0.034$\\
\end{tabular}\end{center}

\Needspace{16\baselineskip}\par\indent A very cursory comparison of my observations and \textsc{Olbers}’s first observations with my elements gave me the following results.
\SourceLeaf{342}
\begin{center}\fontsize{10}{12}\selectfont\setlength{\tabcolsep}{7pt}\renewcommand{\arraystretch}{1.15}\begin{tabular}{lr|r|r}
 & & R.A. & Decl.\\\hline
Aug. & 23 & $+127''$ & $+18''$\\
Sept. & 4 & $+30$ & $-9$\\
 & 6 & $+9$ & $-11$\\
 & 7 & $-54$ & $+20$\\
 & 9 & $-65$ & $-64$\\
 & 14 & $-140$ & .\;.\;.\;.\\
 & 15 & $-195$ & $-76$\\
\end{tabular}\end{center}

\par\indent If the rapidly growing difference in right ascension seen here is real, I would be inclined to regard it already as a sign that the orbit is elliptical.”
